\section{Resultados Gerais}
\label{sec:resultados-gerais}

A Tabela~\ref{tab:resultados-gerais-benchmark} apresenta a síntese dos resultados obtidos pelos algoritmos Morpheus, Sunny e Rocket nos dois cenários avaliados. Foram consideradas as principais métricas do benchmark: quantidade de requisições, taxa de sucesso, latência média, p95, uso médio de memória RAM, pico de memória RAM, número médio de gerações, ocorrências de parada antecipada, erro médio e melhor score.

\begin{table}[H]
\centering
\caption{Resultados gerais do benchmark dos algoritmos Morpheus, Sunny e Rocket}
\label{tab:resultados-gerais-benchmark}
\begingroup
\small
\setlength{\tabcolsep}{4pt}
\renewcommand{\arraystretch}{1.3}

\textbf{Desempenho e confiabilidade}
\vspace{0.15cm}

\begin{tabularx}{\textwidth}{|C{1.6cm}|X|C{2.1cm}|C{1.8cm}|C{2.9cm}|C{2.0cm}|}
\hline
\textbf{Cenário} & \textbf{Algoritmo} & \textbf{Requisições} & \textbf{Sucesso} & \textbf{Latência média (ms)} & \textbf{p95 (ms)} \\
\hline
i3 & Morpheus & 9000 & 99,99\% & 5.343,72 & 6.861,02 \\
\hline
i3 & Sunny & 9000 & 100,00\% & 5.917,48 & 7.546,85 \\
\hline
i3 & Rocket & 9000 & 100,00\% & 803,65 & 986,36 \\
\hline
i5/full & Morpheus & 9000 & 100,00\% & 5.044,99 & 6.058,92 \\
\hline
i5/full & Sunny & 9000 & 100,00\% & 5.832,73 & 7.083,80 \\
\hline
i5/full & Rocket & 9000 & 100,00\% & 640,26 & 777,92 \\
\hline
\end{tabularx}

\vspace{0.3cm}
\textbf{Consumo de memória}
\vspace{0.15cm}

\begin{tabularx}{\textwidth}{|C{2.2cm}|X|C{3.6cm}|C{3.6cm}|}
\hline
\textbf{Cenário} & \textbf{Algoritmo} & \textbf{RAM média (MB)} & \textbf{Pico RAM (MB)} \\
\hline
i3 & Morpheus & 69,86 & 73,75 \\
\hline
i3 & Sunny & 835,63 & 1.329,15 \\
\hline
i3 & Rocket & 13,49 & 24,12 \\
\hline
i5/full & Morpheus & 72,69 & 73,63 \\
\hline
i5/full & Sunny & 778,60 & 1.098,75 \\
\hline
i5/full & Rocket & 11,35 & 20,68 \\
\hline
\end{tabularx}

\vspace{0.3cm}
\textbf{Comportamento e resultado evolutivo}
\vspace{0.15cm}

\begin{tabularx}{\textwidth}{|C{1.6cm}|X|C{2.5cm}|C{2.7cm}|C{2.0cm}|C{2.3cm}|}
\hline
\textbf{Cenário} & \textbf{Algoritmo} & \textbf{Gerações médias} & \textbf{Parada antecipada} & \textbf{Erro médio} & \textbf{Melhor score} \\
\hline
i3 & Morpheus & 191,38 & 8.979 & 3,00 & -29.574 \\
\hline
i3 & Sunny & 945,00 & 0 & 3,00 & -29.640 \\
\hline
i3 & Rocket & 191,52 & 8.980 & 3,00 & -29.574 \\
\hline
i5/full & Morpheus & 191,41 & 8.980 & 3,00 & -29.574 \\
\hline
i5/full & Sunny & 945,00 & 0 & 3,00 & -29.640 \\
\hline
i5/full & Rocket & 191,44 & 8.980 & 3,00 & -29.574 \\
\hline
\end{tabularx}
\endgroup
\fonteautores
\end{table}

Os resultados mostram que o Rocket apresentou a menor latência média nos dois cenários avaliados. No cenário i3, o Rocket obteve média de 803,65 ms, enquanto o Morpheus obteve 5.343,72 ms e o Sunny obteve 5.917,48 ms. No cenário i5/full, o Rocket apresentou média de 640,26 ms, contra 5.044,99 ms do Morpheus e 5.832,73 ms do Sunny. Assim, em ambos os ambientes, o Rocket respondeu de forma mais rápida, enquanto o Morpheus permaneceu mais rápido que o Sunny.

O mesmo comportamento aparece no percentil 95 da latência. No cenário i3, o p95 do Rocket foi de 986,36 ms, enquanto o Morpheus registrou 6.861,02 ms e o Sunny atingiu 7.546,85 ms. No cenário i5/full, o Rocket registrou p95 de 777,92 ms, contra 6.058,92 ms do Morpheus e 7.083,80 ms do Sunny. Isso indica que a vantagem do Rocket não ocorreu apenas na média, mas também nas execuções mais lentas da distribuição analisada.

A diferença mais expressiva entre os algoritmos ocorreu no consumo de memória. No cenário i3, o Rocket utilizou em média 13,49 MB de RAM, o Morpheus utilizou 69,86 MB e o Sunny utilizou 835,63 MB. No cenário i5/full, o Rocket registrou 11,35 MB, o Morpheus registrou 72,69 MB e o Sunny utilizou 778,60 MB. Portanto, nos dois cenários, o Sunny apresentou consumo médio de memória muito superior aos demais, enquanto o Rocket apresentou o menor consumo registrado.

Em relação ao número de gerações executadas, os algoritmos também apresentaram comportamentos bastante distintos. O Morpheus executou, em média, aproximadamente 191 gerações nos dois cenários, resultado muito próximo ao observado no Rocket, que também ficou em torno de 191 gerações. Já o Sunny executou, em média, 945 gerações. Esse resultado indica que Morpheus e Rocket acionaram com frequência o critério de parada antecipada, enquanto o Sunny executou próximo ao limite máximo configurado.

O erro médio final foi igual a 3,00 para os três algoritmos nos dois cenários. Esse resultado impede afirmar que uma das implementações produziu grades completamente livres de inconsistências. Portanto, a diferença entre os algoritmos aparece de forma mais clara nas métricas operacionais, como tempo e memória, do que na eliminação total dos erros da solução final.
